\begin{abstract}
Navigation systems implement established routing objectives --- shortest
distance, reactive travel time, capacity-aware load balancing,
reliability-aware routing --- but commit to one by design rather than by traffic
condition. We ask whether that commitment should be made per context, and
separate two questions usually merged: whether the preferred policy varies with
conditions, and whether exploiting that variation is worth anything.

In \NumRunsMain{} microsimulation runs over \NumContexts{} contexts
(\NumSeeds{} paired seeds per policy) spanning demand, signal timing,
penetration, information lag, disruption and alternative capacity, the
preference does vary: three of four policies are strictly best somewhere by a
seed-resolved margin, and \NumParetoNonSingleton{} of contexts have a
non-singleton Pareto set over journey time, \CO{} and stopped delay --- with the
reliability-aware policy minimising stopped delay in \NumSpecCThreePct{} of
contexts while rarely minimising journey time. Exploiting it is nevertheless
worth \NumHeadroom\,s per vehicle, \NumHeadroomPct{} of the single-best fixed
policy's mean journey time, because that policy's average gain where it wins is
\NumAsymRatio$\times$ its average loss where it does not; selecting the wrong
fixed policy costs \NumFixedPenPctAbsPOne{} of the same baseline.

Within that margin, a selector predicting the winner's identity is \emph{more
accurate and more costly} than the fixed policy, whereas one predicting each
policy's \emph{advantage} closes \NumBFourGapClosed{} of the available gap.
Across \NumOODNTotal{} shift splits the unguarded selector helps on
\NumOODNHelps{}, is indistinguishable on \NumOODNNeutral{} and harms on
\NumOODNHarms{}; a support-and-confidence gate recovered the fixed policy
exactly on every harmful split, at the cost of the in-distribution gain, and one
split shows a shift the gate cannot see because it leaves the feature space
unchanged. The benchmark separates three questions usually treated as one:
complementarity, decision value, and deployability. We introduce no routing,
learning or multi-criteria decision algorithm.

\keywords{Routing-policy selection \and Algorithm selection \and Decision
regret \and Microsimulation \and Distribution shift \and Selective prediction.}
\end{abstract}
